\documentclass[11pt]{article}
\usepackage{amsmath,amssymb}
\usepackage{graphicx}
\usepackage[hidelinks]{hyperref}
% xepersian must be loaded last.
\usepackage{xepersian}
\settextfont[Path=fonts/,Extension=.ttf,UprightFont=*-Regular,BoldFont=*-Bold,ItalicFont=*-Regular,
  ItalicFeatures={FakeSlant=0.2}]{Vazirmatn}
\setlatintextfont[Path=fonts/,Extension=.ttf,UprightFont=*-Regular,BoldFont=*-Bold,ItalicFont=*-Regular,
  ItalicFeatures={FakeSlant=0.2}]{Vazirmatn}

\title{عنوان}
\author{نویسنده}
\date{\today}

\begin{document}
\maketitle

\section{مقدمه}
متن را اینجا بنویسید. واژه‌های انگلیسی را در \lr{English} بگذارید و مقاله‌ها را با دکمه‌ی \lr{Cite} ارجاع دهید.

% plain-fa (persian-bib) sets English entries left to right; give a Persian entry `language = {Persian}`.
\bibliographystyle{plain-fa}
\bibliography{references}
\end{document}
